\documentclass[10pt,a4paper]{amsart}
\usepackage[margin=19mm]{geometry}
\usepackage{amsmath,amssymb,mathtools}
\usepackage[scr=rsfs,cal=euler]{mathalfa}
\usepackage{xfrac}
\usepackage[unicode,hidelinks,bookmarks=false]{hyperref}
\newcommand{\ie}{i.e.\ }
\newcommand{\cf}{cf.\ }
\newcommand{\resp}{respectively\ }
\newcommand{\Subsection}{\subsection}
\providecommand{\nobreakdash}{\nobreak\hbox{-}}
\setlength{\emergencystretch}{2em}
\multlinegap0pt
\usepackage{bm}
\usepackage{bbm}
\usepackage{dsfont}
\usepackage{xspace}
\usepackage{cancel}
\usepackage{mathtools}
\usepackage{amsthm}
\makeatletter
\@ifundefined{lemma}{\newtheorem{lemma}{Lemma}}{}
\@ifundefined{theorem}{\newtheorem{theorem}{Theorem}}{}
\makeatother
\makeatletter
\def\R{{\mathbb{R}}}
\def\calA{{\mathcal{A}}}
\def\calL{{\mathcal{L}}}
\def\la{\lambda}
\@ifundefined{proof}{\newenvironment{proof}{{\it Proof: \enspace}}{\hfill $\blacksquare$\par}}{}
\makeatother
\title[A Spectral-based ISS small-gain theorem for boundary control]{A Spectral-based ISS small-gain theorem for boundary control systems with infinite couplings}
\author{Yassine El Gantouh, Jun Zheng, Guchuan Zhu,, Dingshi Li}
\date{Source: arXiv:2604.11031v1 (CC BY 4.0)}
\begin{document}
\maketitle

\noindent\emph{Source and licence.} A Spectral-based ISS small-gain theorem for boundary control systems with infinite couplings. Version arXiv:2604.11031v1 (CC BY 4.0), distributed under the Creative Commons Attribution 4.0 International licence, as recorded in the arXiv licence metadata for this exact version and shown on the abstract page https://arxiv.org/abs/2604.11031v1. Full rights notice: licenses/arxiv-2604.11031-CC-BY-4.0.txt.\par\smallskip

\noindent\textit{Editorial scope.} Counted unit: the Lemma whose source label is \texttt{technical-lem2} in the article (source0.tex lines 663--668, source label \texttt{technical-lem2}), delivered together with the article's own complete original proof of it (source0.tex lines 669--700, one \texttt{proof} environment of 32 lines). No mathematics is written, repaired or re-derived: statement and proof are the article's own text line for line. Required context retained uncounted: BVCP (lines 457--462); Main1 (lines 481--498); resolvent-operator (lines 549--552); P-regularity (lines 571--576); definition-J' (lines 652--655); extension-identityJ' (lines 658--660). Every definition, display and label of those lines is retained unchanged. Omitted from this package and not redistributed: 1--456, 463--480, 499--548, 553--570, 577--651, 656--657, 661--662, 701--1533. Nothing inside a retained excerpt is omitted or altered: every retained body is delivered exactly as the article's lines are written, so no omission receipt of any kind accompanies this package. Citations: the retained excerpts carry no citation command at all, so no bibliography entry of the article is transcribed and no reference is redistributed. Reference adaptation: none was needed and none was made; every reference inside the delivered unit resolves inside it, so no unresolved reference or citation is produced. Printed slips kept literal: no printed word, symbol, hypothesis or number of the counted statement or of its proof was altered. Layout and class: the article's own class is \texttt{autart} with packages including amsmath, bm, times, mathrsfs, bbm, graphicx, dcolumn, graphics, epsfig, amssymb, amsfonts, fancyhdr, color, subfigure; the wrapper is the lane's pinned \texttt{amsart} editorial wrapper at 10pt on A4, which restates the theorem-like environments the retained text needs and adds the article's own macro definitions that the retained text needs (\texttt{\textbackslash R}, \texttt{\textbackslash calA}, \texttt{\textbackslash calL}, \texttt{\textbackslash la}), with extra wrapper packages (bm, bbm, dsfont, xspace, cancel, mathtools). The delivered numbering is the unit's own. Assets: no retained span contains a figure environment, a graphics-inclusion command, a PDF-page inclusion, a table or a TikZ picture; no image file is copied into this package, the package declares no asset dependency and no image file is redistributed with it. Licence: version arXiv:2604.11031v1 of this article is distributed under the Creative Commons Attribution 4.0 International licence, as recorded in the arXiv licence metadata for this exact version and shown on the abstract page https://arxiv.org/abs/2604.11031v1. A full rights notice is delivered at licenses/arxiv-2604.11031-CC-BY-4.0.txt. Characters: every retained excerpt is pure ASCII, so the delivered engine, XeLaTeX with its default Latin Modern text font, reports no missing character.
    	\begin{align}\label{BVCP}
		\begin{cases}
			\dot{z}(t) = \tilde{A} z(t), & t > 0, \quad z(0)=z_0, \\
			G z(t) = P z(t) + K u(t), & t > 0.
		\end{cases}
	\end{align}

	\begin{theorem}\label{Main1}
		Let $X,Y$, and $U$ be Banach lattices, and let $p \in [1,\infty]$.
		Assume the following conditions hold:
		\begin{itemize}
			\item[\bf{(A1)}] $A := \tilde{A}\vert_{\ker G}$ is a densely defined resolvent positive operator and for some $\lambda_0 > s(A)$ there exists $c > 0$ such that
			\begin{align}\label{inverse2}
				\Vert R(\lambda, A)x \Vert \geq c \Vert x \Vert, \qquad \forall \lambda \ge \lambda_0, \, x \in {X_+}.
			\end{align}
			\item[\bf{(A2)}] $G$ is surjective and $D_{\lambda} \ge 0$  {for all sufficiently large $\lambda\in \R$}.
			\item[\bf{(A3)}] $P$ is positive and $r(P D_{\lambda_1}) < 1$ for some $\lambda_1 > s(A)$.
			\item[\bf{(A4)}] For all $y \in Y$ and $y^* \in Y'$:
			$$\lim_{\R \ni \lambda \to +\infty} \langle P D_\lambda y, y^* \rangle_{Y,Y'} = 0.$$
		\end{itemize}
		Then, the  boundary control system \eqref{BVCP} is well-posed in the sense that for every initial value $z_0 \in X$ and input $u \in L^p_{loc}(\R_+;U)$, there exists a unique mild solution $z \in {C}(\R_+; X)$, and for every $\tau > 0$ there exists a constant $C_{\tau,p} > 0$, which is dependent of $p$ when $p\in(1,\infty)$ and  independent of $p$ when $p=\infty$, such that
		\begin{align*}
			\Vert z(\tau) \Vert_X  \le C_{\tau,p} \left( \Vert z_0 \Vert_X  + \Vert u \Vert_{L^p(0,\tau;U)}   \right).
		\end{align*}
	\end{theorem}

 	\begin{align}\label{resolvent-operator}
		R(\lambda, \calA) = (\mathcal{I}-D_\lambda P)^{-1} R(\lambda, A)
		= R(\lambda, A) + D_\lambda (\mathcal{I}-P D_\lambda)^{-1} P R(\lambda, A).
	\end{align}

	\begin{lemma}\label{P-regularity}
		Let $A := \tilde{A}|_{\ker G}$ be the generator of a positive $C_0$-semigroup $T$ on $X$ and $P \ge 0$. Suppose assumptions {\bf (A2)} and {\bf (A4)} are satisfied. Then
		\begin{align}\label{Limit-Px}
			\lim_{\la \to +\infty} \Vert \la P R(\la,A)x-Px\Vert_Y=0, \qquad \forall x\in Z.
		\end{align}
	\end{lemma}

	\begin{align}
		 J^\calA x:=&\lim_{\la \to +\infty} \Vert \la R(\la,\calA)x\Vert_{-1,A},\label{definition-J'}\\
		  {\bm D}(J^\calA):=&\{x\in X_{-1,\calA}: \text{ the limit in \eqref{definition-J'}}\text{ exists in } X_{-1,A}\}. \notag
	\end{align}

	\begin{equation}\label{extension-identityJ'}
		J^\calA x=x,\qquad \forall x\in X.
	\end{equation}

	\begin{lemma}\label{technical-lem2}
		Consider the setting of Theorem \ref{Main1}. Then, $Z_{-1,\calA}\subset {\bm D}(J^\calA)$ and
		\begin{align}\label{boundness-J'}
			J^\calA \in \calL(Z_{-1,\calA},Z_{-1,A}).
		\end{align}
	\end{lemma}

	\begin{proof}
		First, we establish that for all $x\in {\bm D}(\tilde{A})$,
		\begin{align}\label{regularity-P-wrt-calA}
			\lim_{\la \to +\infty}\la PR(\la,\calA)x=Px.
		\end{align}
		Let $x\in {\bm D}(\tilde{A})$. Using \eqref{resolvent-operator},  Lemma \ref{P-regularity}, and assumption {\bf (A4)}, we obtain
		\begin{equation*}
			\lim_{\la \to +\infty}\la PR(\la,\calA)x =\lim_{\la \to +\infty}\la PR(\la,A)x+\lim_{\la \to +\infty}  PD_\la (\mathcal{I}-PD_\la )^{-1}\la PR(\la,A)x
			=Px.
		\end{equation*}
		Now, let $z\in Z_{-1,\calA}$. By \eqref{definition-J'}, we must show that $\la R(\la,\calA)z$ converges in $X_{-1,A}$ as $\la \to +\infty$. Fix $ s>s(\calA)$. Applying \eqref{resolvent-operator}, we obtain
		\begin{align*}
			R(s,A)\la R(\la,\calA)z&=(R(s,\calA)-D_s PR(s,\calA))\la R(\la,\calA) z\\
			&=\la R(\la,\calA) R(s,\calA)z-D_s P\la R(\la,\calA)R(s,\calA)z.
		\end{align*}
		Taking the limit $\lambda \to +\infty$ and invoking \eqref{extension-identityJ'} and \eqref{regularity-P-wrt-calA}, we obtain
		\begin{align*}
			\lim_{\la \to +\infty}R(s,A)\la R(\la,\calA)z&=\lim_{\la \to +\infty}\la R(\la,\calA) R(s,\calA)z-D_s \lim_{\la \to +\infty}\la P R(\la,A)R(s,\calA)z\\
			&=R(s,\calA)z-D_s PR(s,\calA)z.
		\end{align*}
		This proves $Z_{-1,\calA}\subset {\bm D}(J^\calA)$  with
		\begin{align*}
			R(s,A)J^\calA z=R(s,\calA)z-D_s PR(s,\calA)z.
		\end{align*}
		From the above expression, we derive the following norm estimate
		%{which implies that}
		\begin{align*}
			\Vert J^\calA z\Vert_{Z_{-1,A}}=\Vert R(s, A)J^\calA z\Vert_{Z}
			\le \left(1+\Vert D_s P\Vert_{\calL(Z,Y)}\right)\Vert z\Vert_{Z_{-1,\calA}},
		\end{align*}
		for all $z\in Z_{-1,\calA}$. Thus, $J^\calA \in \calL(Z_{-1,\calA},Z_{-1,A})$. This ends the proof.
	\end{proof}

\end{document}
